\begin{lemma}[Infimum Approximation Lemma]
Let $A \subseteq \mathbb{R}$ be nonempty and bounded below, and let
\[
i = \inf A .
\]
Then for every $\varepsilon > 0$ there exists $a \in A$ such that
\[
a < i + \varepsilon .
\]
\end{lemma}

\begin{proof}
Assume for contradiction that there exists $\varepsilon_0 > 0$ such that

\[
\forall a \in A,\quad i + \varepsilon_0 \le a .
\]

This means that

\[
i + \varepsilon_0
\]

is a lower bound for $A$.

But $i = \inf A$ is the \textbf{greatest lower bound}. However,

\[
i < i + \varepsilon_0 .
\]

Thus $i + \varepsilon_0$ would be a lower bound greater than $i$,
which contradicts the fact that $i$ is the infimum of $A$.

Therefore our assumption is false. Hence there exists $a \in A$
such that

\[
a < i + \varepsilon .
\]
\end{proof}

\begin{lemma*}[Infimum Approximation Lemma]
Let $A \subseteq \mathbb{R}$ be nonempty and bounded below, and let
\[
i = \inf A .
\]
Then for every $\varepsilon > 0$ there exists $a \in A$ such that
\[
a < i + \varepsilon .
\]
\end{lemma*}

\begin{proof}
Assume for contradiction that there exists $\varepsilon_0 > 0$ such that

\[
\forall a \in A,\quad i + \varepsilon_0 \le a .
\]

This means that

\[
i + \varepsilon_0
\]

is a lower bound for $A$.

But $i = \inf A$ is the \textbf{greatest lower bound}. However,

\[
i < i + \varepsilon_0 .
\]

Thus $i + \varepsilon_0$ would be a lower bound greater than $i$,
which contradicts the fact that $i$ is the infimum of $A$.

Therefore our assumption is false. Hence there exists $a \in A$
such that

\[
a < i + \varepsilon .
\]
\end{proof}